% Standalone section, ready to include in the parent manuscript.
% Requires amsmath, amssymb, amsthm, hyperref; theorem/lemma/proposition/remark
% environments may be supplied by the parent.
\section{An exact rectangular backend from the new characteristic-zero bounds}
\label{sec:number-field-backend}

The matrix-multiplication input is updated at this pinned OpenAI Math commit:
\begin{center}\small
\texttt{7ef5c0cc5ca4b0e072b7a050fca09ab24bc6fa3a}.
\end{center}
The October~2 preprint \cite{matrix94} \emph{An Upper Bound of $9/4$ for the Matrix
Multiplication Exponent} proves
\[
 \omega_{\mathbb C}\leq\frac94.
\]
Its proof bounds the exact-rank exponent, then selects one fixed finite
exact rank scheme for every positive accuracy. Its final remark explicitly
places the coefficients in a number field by the Nullstellensatz.
The September~24 preprint \cite{matrixdual} \emph{Complex Matrix Multiplication Below
$2.258$ and Rectangular Bounds} proves
\[
 \alpha_{\mathbb C}>0.465=\frac{93}{200},
\]
and its characteristic-zero dual-bound corollary transfers this statement
to every characteristic-zero field. We use these two primary results as
mathematical dependencies; their entire tensor constructions are not
reproved here. The additional argument below supplies the exact
bit-complexity and residue-ring interface that the arithmetic statements
alone do not provide.

Put $\rho=93/200$. Fix a rational accuracy $\varepsilon>0$, independent
of the SAT instance. In the numerical certificates we take
$\varepsilon=10^{-5}$. Every constant in this section may depend on this
one fixed accuracy and on the resulting finite schemes, but not on matrix
dimensions, entry bit lengths, the SAT input, or a recursion depth.

\subsection{A finite exact scheme is effective}

\begin{lemma}[Effective algebraic schemes]
\label{lem:effective-algebraic-schemes}
For every rational $\delta>0$, one can deterministically find a number
field $K$ and two fixed exact bilinear schemes over $K$ as follows:
\begin{enumerate}
\item a square scheme for $u\times u$ multiplication, with $u\geq2$
and rank budget $r_{\rm sq}<u^{9/4+\delta}$;
\item a rectangular scheme for the dimensions $(a,b,a)$, with
$a\geq2$, $1\leq b\leq a$, $b^{200}\geq a^{93}$, and rank budget
$r_{\rm fl}<a^{2+\delta}$.
\end{enumerate}
The search is a finite preprocessing computation depending only on
$\delta$. No assertion of a practical search time or small scheme is made.
\end{lemma}

\begin{proof}
For the square scheme, the exact-rank exponent in the October~2 theorem
is at most $9/4$, so some fixed $u$ and exact rank budget satisfy the
strict displayed inequality.

For the rectangular scheme, choose $k>\rho$ with
$\omega_{\mathbb C}(1,k,1)=2$. Such a $k$ exists by the strict dual bound.
The exact-rank limit for rectangular matrix multiplication, proved in
the September~24 paper, supplies arbitrarily large fixed values $t$ and
exact schemes of dimensions
\[
 a=\lceil e^t\rceil,\qquad b=\lceil e^{kt}\rceil,
 \qquad \log r_{\rm fl}<(2+\delta/2)t.
\]
One may take $k<1$, by padding if necessary. For sufficiently large $t$,
$b\leq a$ and $b^{200}\geq a^{93}$, since $k>93/200$.
Also $r_{\rm fl}<a^{2+\delta}$. Thus a finite scheme with the required
integer inequalities exists.

For either specified dimension triple and rank budget, express each
rank-one term as a product of three linear forms. Equating the
coefficients with those of the matrix-multiplication tensor gives
finitely many polynomial equations over $\mathbb Q$:
\[
 \sum_{\nu=1}^{r} A_{\nu,i}B_{\nu,j}C_{\nu,k}
       =T_{ijk}\in\{0,1\}.
\]
A complex solution makes this ideal proper. The Nullstellensatz gives
a solution over $\overline{\mathbb Q}$, whose finitely many coefficients
lie in a number field. Taking the compositum puts both schemes in one
number field $K$.

This existence assertion is effective in the elementary, unbounded-search
sense needed here. Enumerate finite dimension triples, rank budgets,
number fields specified by irreducible rational polynomials, and rational
coordinate descriptions of all scheme coefficients. Test the coefficient
identities exactly in the proposed field and test the rank/shape
inequalities by integer powers, since $\delta$ and $\rho$ are rational.
Enumeration by increasing finite description length visits every finite
candidate. Exact field arithmetic and irreducibility testing over
$\mathbb Q$ are effective. An algebraic solution described above will
eventually be visited, so the search terminates. One can instead use
effective polynomial-system elimination, but no bound for that elimination
is needed. After this preprocessing, the two schemes are fixed finite
data in a single algorithm for all input sizes.
\end{proof}

\subsection{Coefficient growth in the whole recursive algorithm}

\begin{lemma}[Bit cost of a fixed algebraic bilinear recursion]
\label{lem:height-bridge}
Fix finitely many exact bilinear matrix schemes with coefficients in a
number field $K$, and use them by recursive blocking, tensor composition,
permutation of the three tensor legs, zero padding, and ordinary blocking.
Suppose every recursion path has length $h$, and let $N$ be the largest matrix dimension.
For integer input entries of at most $B$ bits, all intermediate number-field
coordinates can be represented exactly by rational numbers with
$O(B+h+\log N+1)$ bits in numerator and denominator.
Consequently each scalar operation of the whole algebraic computation
costs $(B+h+\log N+1)^{O(1)}$ bit operations. Simulating in a fixed rational basis
changes the whole arithmetic operation count by a constant factor; it does
not multiply the rank budget by the field degree at every recursion level.
\end{lemma}

\begin{proof}
Choose a rational basis $e_1=1,e_2,\ldots,e_d$ of $K$. The degree $d$ is
fixed. Write every scheme coefficient in this basis and write the fixed
structure constants in
\[
 e_i e_j=\sum_{k=1}^{d}\mu_{ijk}e_k.
\]
Choose an integer $D\geq2$ clearing their finitely many rational
denominators. Then
\[
 \mathcal A=\bigoplus_{i=1}^{d}\mathbb Z[1/D]e_i
\]
is a subalgebra of $K$ containing all coefficients. There is a fixed
constant $C$ bounding all relevant coordinate magnitudes, the number of
terms in each fixed linear form, and the number of coordinates used in
one product.

In a recursive bilinear scheme, the operands passed down any branch are
linear combinations of the original left inputs and, separately, of the
original right inputs. At one level each coordinate is a fixed linear
combination of a bounded number of coordinates. After $j$ levels its
denominator divides $D^{c j}$, and its absolute value is at most
$C^{c j}2^B$, for a fixed $c$. Thus it has $O(B+j+1)$ bits in rational
coordinates. This conclusion follows equally by induction with a common
denominator; reduced fractions are not required.

At a leaf, one multiplies a left linear form by a right linear form.
The two denominator exponents add and the coordinate structure constants
add only a fixed amount. The resulting coordinate numerators have
$O(2B+h+1)$ bits. On the upward path only fixed linear combinations of
these bilinear values are taken. Each level increases the logarithmic
magnitude and common-denominator exponent by at most a fixed constant.
Hence the final bound is $O(B+h+\log N+1)$ throughout.

This argument uses the bilinear structure of the recursion: no sequence
of repeated squarings of an arbitrary computed value is present.
Equivalently, each leaf product has degree one in each original operand,
and every later operation is linear in such products. The exponentially
large expanded rank decomposition need not be stored or constructed.
Ordinary block summation adds at most $O(\log N)$ to the bit length, because there are polynomially many summands in $N$. The ordinary recursive implementation already performs the operations
whose heights were bounded.

A number-field addition or multiplication uses a fixed number of rational
coordinate operations, because $d$ and its multiplication table are fixed.
Elementary integer arithmetic and exact rational normalization on
$O(B+h+\log N+1)$ bits take polynomial time. This simulates each node of the
already chosen recursion over $K$. It introduces a constant factor per
executed operation, rather than replacing a rank-$r$ recurrence by a
rank-$(d r)$ recurrence.
\end{proof}

In particular $h=O(1+\log N)$ for the fixed schemes on matrices of maximum
dimension $N$. A depth-first bilinear implementation forms one child
product at a time, adds it into its parent output, and releases it.
At each ancestor it retains only a fixed number of input/output-sized
arrays. Their dimensions decrease geometrically down a nontrivial
recursion path. Its bit space is therefore a polynomial factor in
$B+\log N+1$ times the largest input/output array, with a harmless
additional logarithmic factor if an implementation retains all active
descriptors.

\subsection{Integer outputs and residue rings, including characteristic two}

\begin{proposition}[Exact lifting to the rings used by the contractor]
\label{prop:residue-ring-bridge}
Let $p,\ell,t$ and the bit lengths of all ring data be polynomially
bounded in a parameter $s$. Let $f\in\mathbb Z[U]$ be monic of degree
$t$, with polynomially bounded coefficient bit lengths, and put
\[
 R=(\mathbb Z/2^p\mathbb Z)[U,V]/(f(U),V^\ell).
\]
For matrices over $R$ whose dimension logarithms are polynomially bounded
in $s$, any backend furnished by Lemma~\ref{lem:effective-algebraic-schemes}
can be evaluated exactly with only a factor polynomial in $s$ overhead
over its matrix arithmetic count. The factor's exponent and all scheme
data are fixed independently of $s$. No denominator is inverted modulo
$2^p$.
\end{proposition}

\begin{proof}
First consider integer matrices. Run the entire fixed recursive algorithm
over the number field $K$, represented in the rational basis above.
The bilinear identities hold exactly, so on rational inputs the output
lies in the one-dimensional rational subspace $\mathbb Q e_1$ and equals
the ordinary product. On integer inputs it is an integer. All other
coordinates are exactly zero. Project to the coordinate of $e_1$ and
perform its final exact rational division. Its denominator divides
$D^{O(\log N)}$, whose bit length is $O(\log N)$; its numerator and
denominator have the bit bounds of Lemma~\ref{lem:height-bridge}.
An exact integer quotient is obtained, not a numerical approximation.
The integer $D$ may be even. Nothing has been reduced modulo $2^p$
during this rational computation, so even denominators cause no problem.

For entries in $R$, take the canonical polynomial representatives
\[
 \sum_{0\leq i<t,\,0\leq j<\ell} c_{ij}U^iV^j,
 \qquad 0\leq c_{ij}<2^p.
\]
Each coordinate is an integer matrix. Multiply every required pair of
coordinate matrices by the exact integer procedure, collect equal
monomials, discard powers $V^j$ with $j\geq\ell$, and reduce the
$U$-degree using the monic polynomial $f$. Finally reduce integer
coefficients modulo $2^p$.
There are at most $(t\ell)^2$ coordinate matrix products, and all
remaining operations are polynomially many array additions and scalar
integer operations per entry. Monicity makes polynomial reduction exact
over the integers, without division by a nonunit. The coefficient growth
of this finite reduction is polynomial in $s$, since both its degree and
the coefficient bit lengths of $f$ are polynomially bounded.

The natural map $\mathbb Z[U,V]\to R$ is a ring homomorphism, so this
procedure gives the required product in $R$. We do not assert that
canonical lifting from $R$ to integer polynomials is itself a ring
homomorphism. We only lift the given inputs, compute their polynomial
product exactly, and then apply the quotient map. This also leaves the
contractor's immutable creation-precision convention intact.

Here input bit lengths and $\log N$ are polynomial in $s$, so
Lemma~\ref{lem:height-bridge} and the polynomial number of coordinate
products prove the claimed overhead. CRT is optional and unnecessary:
the exact rational/integer route already handles characteristic two,
including the case in which two divides every natural denominator of
the chosen algebraic schemes.
\end{proof}

\subsection{The resulting nine-support rectangular bound}

\begin{theorem}[Updated exact rectangular backend]
\label{thm:updated-rectangle}
Let $x\geq y\geq z\geq0$ be the base-two logarithms of three positive
integer matrix dimensions, in decreasing order. Under the bit-size
conditions of Proposition~\ref{prop:residue-ring-bridge}, the product can
be computed exactly in a polynomial factor times $2^{F_\varepsilon(x,y,z)}$
bit operations, where
\[
 F_\varepsilon(x,y,z)=
 \max\left\{
 (1+\varepsilon/2)(x+y),\quad
 x+\left(\frac{335}{428}+\varepsilon\right)y+
       \frac{50}{107}z
 \right\}.
\]
Its bit space is bounded by a polynomial factor times $2^{x+y}$,
and hence also by a polynomial factor times $2^{F_\varepsilon(x,y,z)}$.
For $\varepsilon=10^{-5}$, this is exactly the maximum of the nine linear
supports obtained by permuting
\[
 (1+\varepsilon/2,1+\varepsilon/2,0),\qquad
 \left(1,\frac{335}{428}+\varepsilon,\frac{50}{107}\right).
\]
\end{theorem}

\begin{proof}
Choose the two schemes with a smaller fixed accuracy, for instance
$\delta=\varepsilon/4$, so that logarithmic recursion overhead can be
absorbed into the stated $\varepsilon$ slack. Their coefficient fields
are combined once, as above. First block the largest dimension down to
the scale of the middle dimension, paying $2^{x-y+O(1)}$ blocks.

Suppose $z\geq\rho y$. Define
\[
 t_0=\frac{y-z}{1-\rho},\qquad
 t_1=\frac{z-\rho y}{1-\rho}.
\]
These are nonnegative, satisfy $t_0+t_1=y$, and decompose the remaining
shape into a flat shape $(t_0,\rho t_0,t_0)$ and a square shape of
side exponent $t_1$. Tensor the corresponding fixed recursive schemes.
Rounding their recursion depths upward changes their large dimensions
by fixed factors. The flat scheme may have $b>a^\rho$; this only pads
the smaller dimension and it stays no larger than the large dimensions,
because $b\leq a$. Thus the largest padded array still has
$2^{2y+O(1)}$ entries per large-dimension block.

For clarity, the rounded depths can be chosen without an oracle for
logarithms. If the middle and small dimensions are $n_y,n_z$, choose the
least nonnegative $h_0,h_1$ satisfying respectively
\[
 a^{107h_0}n_z^{200}\geq n_y^{200},\qquad
 u^{107h_1}n_y^{93}\geq n_z^{200}.
\]
These are exactly the upward roundings of $t_0/\log_2 a$ and
$t_1/\log_2 u$. Fixed-degree integer powers and comparisons determine
them in polynomial bit time. They guarantee the required dimensions;
the two rounding errors are bounded constants.

The fixed-scheme cost, with its logarithmic recursion factors absorbed
by the reserved accuracy, is at most
\begin{align*}
 x-y+(2+\varepsilon)t_0+(9/4+\varepsilon)t_1
 &=x+\left(1-\rho\frac{1/4}{1-\rho}
                  +\varepsilon\right)y+
                  \frac{1/4}{1-\rho}z\\
 &=x+\left(\frac{335}{428}+\varepsilon\right)y+
                  \frac{50}{107}z.
\end{align*}
Here an additive constant in the exponent is absorbed in a fixed
multiplicative constant. If $z<\rho y$, pad to the flat shape. Its
exponent is at most $x+y+\varepsilon y$, which is bounded by
$(1+\varepsilon/2)(x+y)$ because $x\geq y$. These two cases prove the
time bound. The largest input/output array has logarithmic size $x+y$;
the depth-first implementation and the integer-ring lifting above give
the space bound.

The bilinear tensor decomposition can be executed with any of its three
legs designated as the output. Permuting or transposing the three legs
therefore supplies every ordering of the dimension triple, with the
input/output arrays still charged. Finally,
\[
 1>\frac{335}{428}+10^{-5}>\frac{50}{107}>0.
\]
The rearrangement inequality says that the largest permutation of the
second support assigns these three weights to $x,y,z$ in decreasing
order; the first support has three distinct permutations. Thus the
support maximum is exactly $F_\varepsilon$, not an uncharged choice of
an arithmetic algorithm. In particular the full-array branch is retained.
\end{proof}

\paragraph{What is uniform and what remains asymptotic.}
The accuracy, schemes, number field, basis, and rational multiplication
table are fixed before any SAT input. Their effective search can be
performed once or included as a finite initialization. Matrix sizes and
entry bit lengths vary; their handling is uniform and deterministic.
The proof establishes an eventual complexity bound, not a supplied fast
finite matrix implementation or a practical crossover estimate. In
particular the enormous unknown fixed search and scheme constants must
not be described as negligible in measured implementations. They are
legitimate fixed constants in the asymptotic theorem.

\paragraph{Pinned primary sources.}
The source files and their blob hashes are:
\begin{enumerate}
\item October~2 square theorem and exact-rank/algebraic-coefficient
remark: \href{https://github.com/openai/math/blob/7ef5c0cc5ca4b0e072b7a050fca09ab24bc6fa3a/preprints/Matrix-Multiplication-Nine-Fourths-October-2-2026/build/paper.tex}{\emph{An Upper Bound of $9/4$ for the Matrix Multiplication Exponent}};
blob \texttt{4492e66c707b9329f763c46e7cc113ff21d2edc0}.
\item September~24 strict dual theorem:
\href{https://github.com/openai/math/blob/7ef5c0cc5ca4b0e072b7a050fca09ab24bc6fa3a/preprints/Complex-Matrix-Multiplication-Below-2.258-and-Rectangular-Bounds-September-24-2026/build/introduction.tex}{\emph{Complex Matrix Multiplication Below $2.258$ and Rectangular Bounds}, introduction};
blob \texttt{6b5dd1a6ac7f36b8dc046833960735094ca4e73f}.
\item Same paper, coefficient-equation descent and characteristic-zero
dual-bound corollary:
\href{https://github.com/openai/math/blob/7ef5c0cc5ca4b0e072b7a050fca09ab24bc6fa3a/preprints/Complex-Matrix-Multiplication-Below-2.258-and-Rectangular-Bounds-September-24-2026/build/rectangular/08-comparison.tex}{field-transfer section};
blob \texttt{b7fa963d5e2659639b047e3a7dc179a090acd40a}.
\end{enumerate}
